\section*{結言}
今回の観測では、大規模モデルの公開と限定提供が対照的な動きとして表れた。Mistral Large 4はオープンウェイトでの利用可能性が話題となり、Gemini 4 Carbonはサイバー防御向けの性能とアクセス制御が議論された。公開形態が異なっても、実際の利用条件、性能の再現性、運用コストを確認する必要がある点は共通する。

Project Suncatcherの宇宙TPU実験とGenie 4の研究プレビューは、AIを実行する環境やシミュレーション能力の拡張を示す。これに対し、モデル監視、科学研究支援、音声エージェント、デザイン編集の話題は、技術を具体的な用途へ接続する段階に関心が移っていることを示した。カナダの出資報道やRedditの検索流入減少は、政策と事業構造への影響も浮かび上がらせた。

一方、性能順位、商用ライセンス、利用上限変更、買収条件、オープンソース公開範囲には未確認事項が残る。観測窓より前の元発表と窓内の反応も混在している。次の焦点は、各社の一次資料による条件の明確化と、実利用・独立再現に基づく評価である。
